
% Options for packages loaded elsewhere
\PassOptionsToPackage{unicode}{hyperref}
\PassOptionsToPackage{hyphens}{url}
\documentclass[draft]{report}
\usepackage{xcolor}
\usepackage{amsmath,amssymb}
\setcounter{secnumdepth}{5}
\usepackage{iftex}
\usepackage[T1]{fontenc}
\usepackage[utf8]{inputenc}
\usepackage{textcomp} % provide euro and other symbols
\usepackage[]{fontawesome}

% Use upquote if available, for straight quotes in verbatim environments
\IfFileExists{upquote.sty}{\usepackage{upquote}}{}
\IfFileExists{microtype.sty}{% use microtype if available
  \usepackage[]{microtype}
  \UseMicrotypeSet[protrusion]{basicmath} % disable protrusion for tt fonts
}{}
\makeatletter
\@ifundefined{KOMAClassName}{% if non-KOMA class
  \IfFileExists{parskip.sty}{%
    \usepackage{parskip}
  }{% else
    \setlength{\parindent}{0pt}
    \setlength{\parskip}{6pt plus 2pt minus 1pt}}
}{% if KOMA class
  \KOMAoptions{parskip=half}}
\makeatother
\setlength{\emergencystretch}{3em} % prevent overfull lines
\providecommand{\tightlist}{%
  \setlength{\itemsep}{0pt}\setlength{\parskip}{0pt}}
\usepackage[explicit]{titlesec}
\usepackage[svgnames]{xcolor}
\usepackage{enumitem}
\usepackage{multirow}
\usepackage{array}
\usepackage{calc}
\usepackage{layouts}

% section headers
\titlespacing{\chapter}{0pt}{0pt}{0pt}

\titleformat{\chapter}
  [hang]
  {\normalfont\sffamily\Large\color{white}}
  {}
  {0.5ex}
  {\colorbox{Navy}{\makebox[4em][l]{Article \thechapter} \makebox[\textwidth - 5em][l]{#1} }}  
  %{\colorbox{RoyalBlue}{%
  %  \begin{tabular}{ p{11em} p{\textwidth - 10em} }
  %    Article \thechapter & #1 \\
  %  \end{tabular}
  %}}

% \newcommand{\chapterbreak}{\addpenalty{-500}\vspace*{0pt}}

\titleformat{\section}
  {\normalfont\sffamily}
  {}
  {0.5ex}
  {#1}

% enumerations
\renewcommand{\labelenumi}{\arabic{enumi}}
\renewcommand{\labelenumii}{\arabic{enumi}.\arabic{enumii}}
\renewcommand{\labelenumiii}{\arabic{enumi}.\arabic{enumii}.\arabic{enumiii}}
\renewcommand{\labelenumiv}{\arabic{enumi}.\arabic{enumii}.\arabic{enumiii}.\arabic{enumiv}}

\setlist[enumerate,1]{leftmargin=3ex, rightmargin=0pt, topsep=0pt,%
                       labelsep=1ex, listparindent=0pt, itemsep=1ex}
\setlist[enumerate,2]{leftmargin=4ex, rightmargin=0pt, labelsep=1ex, labelwidth=0em}
\setlist[enumerate,3]{leftmargin=6ex, rightmargin=0pt, labelsep=1ex, labelwidth=0em}
% \setlength{\listparindent=0pt},

% environment for the explanations
\newenvironment{expl}[0]%
  {\rule{\textwidth}{.2pt}\footnotesize}%
  {\rule{\textwidth}{.2pt}\normalsize}



\usepackage{bookmark}
\IfFileExists{xurl.sty}{\usepackage{xurl}}{} % add URL line breaks if available
\urlstyle{same}
\hypersetup{
  hidelinks,
  pdfcreator={LaTeX via pandoc}}

\author{}
\date{}


\begin{document}

\chapter{Purpose}\label{purpose}

\begin{enumerate}
\tightlist
\item
\input{1.1.txt}
\end{enumerate}

\begin{expl}
\section{(Interpretation)}\label{interpretation}
\input{1.i_1.txt}
\end{expl}

\chapter{Application}\label{application}

\begin{enumerate}
\item
\input{2.1.txt}
\item
\input{2.2.txt}
\item
\input{2.3.txt}
\end{enumerate}

\begin{expl}
\section{(Explanation\#1)}\label{explanation1}
\input{2.e_1.txt}
\end{expl}


\chapter{Competition Area}\label{competition-area}

\begin{enumerate}
\item
\input{3.1.txt}
  \begin{enumerate}
  \item
  \input{3.1.1.txt}
  \item
  \input{3.1.2.txt}
  \end{enumerate}
\item
\input{3.2.txt}
  \begin{enumerate}
  \item
  \input{3.2.1.txt}
  \item
  \input{3.2.2.txt}
  \item
  \input{3.2.3.txt}
  \item
  \input{3.2.4.txt}
  \item
  \input{3.2.5.txt}
  \item
  \input{3.2.6.txt}
  \item
  \input{3.2.7.txt}
  \end{enumerate}
\item
\input{3.3.txt}
  \begin{itemize}
  \tightlist
  \item
  \input{3.3.1.txt}
  \item
  \input{3.3.2.txt}
  \item
  \input{3.3.3.txt}
  \item
  \input{3.3.4.txt}
  \end{itemize}
\end{enumerate}

\begin{expl}
\section{(Explanation \#1)}\label{explanation-1}
\input{3.e_1.txt}

\section{(Explanation \#2)}\label{explanation-2}
\input{3.e_2.txt}

\section{(Explanation \#3)}\label{explanation-3}
\input{3.e_3.txt}
\end{expl}

%todo: Ide kell a rajz! és landscape oldalon




\chapter{Contestant}\label{contestant}

\begin{enumerate}
\item
\input{4.1.txt}

  \begin{enumerate}
  \item
  \input{4.1.1.txt}
  \item
  \input{4.1.2.txt}
  \item
  \input{4.1.3.txt}
  \item
  \input{4.1.4.txt}
  \item
  \input{4.1.5.txt}
  \end{enumerate}
\end{enumerate}

\begin{expl}
\section{(Interpretation)}\label{interpretation-1}
\input{4.i_1.txt}

\section{(Interpretation)}\label{interpretation-2}
\input{4.i_2.txt}

\section{(Interpretation)}\label{interpretation-3}
\input{4.i_3.txt}
\end{expl}

\begin{enumerate}[resume*]
\item
\input{4.2.txt}

  \begin{enumerate}
  \item
  \input{4.2.1.txt}

    \begin{enumerate}
    \item
    \input{4.2.1.1.txt}
    \end{enumerate}

  \item
  \input{4.2.2.txt}
  \item
  \input{4.2.3.txt}
  \item
  \input{4.2.4.txt}
  \item
  \input{4.2.5.txt}

    \begin{enumerate}
    \item
    \input{4.2.5.1.txt}

    \item
    \input{4.2.5.2.txt}
    \item
    \input{4.2.5.3.txt}
    \end{enumerate}

  \end{enumerate}

\item
\input{4.3.txt}

  \begin{enumerate}
  \item
  \input{4.3.1.txt}
  \item
  \input{4.3.2.txt}
  \item
  \input{4.3.3.txt}
  \item
  \input{4.3.4.txt}
  \end{enumerate}

\end{enumerate}

\begin{expl}
\section{(Explanation \#1)}\label{explanation-1-1}
\input{4.i_1.txt}

\section{(Explanation \#2)}\label{explanation-2-1}
\input{4.i_2.txt}

\section{(Explanation \#3)}\label{explanation-3-1}
\input{4.e_3.txt}

\section{(Explanation \#4)}\label{explanation-4}
\input{4.e_4.txt}

\section{(Explanation \#5)}\label{explanation-5}
\input{4.e_5.txt}

\section{(Explanation \#6)}\label{explanation-6}
\input{4.e_6.txt}

\section{(Explanation \#7)}\label{explanation-7}
\input{4.e_7.txt}

\section{(Explanation \#8)}\label{explanation-8}
\input{4.e_8.txt}

\section{(Explanation \#9)}\label{explanation-9}
\input{4.e_9.txt}

\section{(Explanation \#10)}\label{explanation-10}
\input{4.e_10.txt}
\end{expl}


\chapter{Weight category}\label{weight-category}

\begin{enumerate}

\item
\input{5.1.txt}

\begin{center}
\begin{tabular}{ |l|l|l|l| } 
 \hline
 \multicolumn{2}{|l|}{Men's division} & \multicolumn{2}{|l|}{Women's division}  \\
 \hline
 Under 54 kg & Not exceeding 54 kg & under 46 kg & Not exceeding 46 kg \\
 \hline
 Under 58 kg & Over 54 kg \& Not exceeding 58 kg & Under 49 kg & Over 46 kg \& Not exceeding 49 kg \\
 \hline
 Under 63 kg & Over 58 kg \& Not exceeding 63 kg & Under 53 kg & Over 49 kg \& Not exceeding 53 kg \\
 \hline
 Under 68 kg & Over 63 kg \& Not exceeding 68 kg & Under 57 kg & Over 53 kg \& Not exceeding 57 kg \\
 \hline
 Under 74 kg & Over 68 kg \& Not exceeding 74 kg & Under 62 kg & Over 57 kg \& Not exceeding 62 kg \\
 \hline
 Under 80 kg & Over 74 kg \& Not exceeding 80 kg & Under 67 kg & Over 62 kg \& Not exceeding 67 kg \\
 \hline
 Under 87 kg & Over 80 kg \& Not exceeding 87 kg & Under 73 kg & Over 67 kg \& Not exceeding 73 kg \\
 \hline
 Over 87 kg & Over 54 kg \& Not exceeding 58 kg & Over 73 kg & Over 73 kg \\
 \hline
\end{tabular}
\end{center}

\item
\input{5.2.txt}

\begin{center}
\begin{tabular}{ |l|l|l|l| } 
 \hline
 \multicolumn{2}{|l|}{Men's division} & \multicolumn{2}{|l|}{Women's division}  \\
 \hline
 Under 58 kg & Not exceeding 58 kg & under 49 kg & Not exceeding 49 kg \\
 \hline
 Under 68 kg & Over 58 kg \& Not exceeding 68 kg & Under 57 kg & Over 49 kg \& Not exceeding 57 kg \\
 \hline
 Under 80 kg & Over 68 kg \& Not exceeding 80 kg & Under 67 kg & Over 57 kg \& Not exceeding 67 kg \\
 \hline
 Over 80 kg & Over 50 kg & Over 67 kg & Over 67 kg \\
 \hline
\end{tabular}
\end{center}

\begin{samepage}
\item
\input{5.3.txt}

\begin{center}
\begin{tabular}{ |l|l|l|l| } 
 \hline
 \multicolumn{2}{|l|}{Men's division} & \multicolumn{2}{|l|}{Women's division}  \\
 \hline
 Under 45 kg & Not exceeding 45 kg & under 42 kg & Not exceeding 42 kg \\
 \hline
 Under 48 kg & Over 45 kg \& Not exceeding 44 kg & Under 49 kg & Over 42 kg \& Not exceeding 44 kg \\
 \hline
 Under 51 kg & Over 48 kg \& Not exceeding 46 kg & Under 53 kg & Over 44 kg \& Not exceeding 46 kg \\
 \hline
 Under 55 kg & Over 51 kg \& Not exceeding 49 kg & Under 49 kg & Over 46 kg \& Not exceeding 49 kg \\
 \hline
 Under 59 kg & Over 55 kg \& Not exceeding 52 kg & Under 53 kg & Over 49 kg \& Not exceeding 52 kg \\
 \hline
 Under 63 kg & Over 59 kg \& Not exceeding 55 kg & Under 57 kg & Over 52 kg \& Not exceeding 55 kg \\
 \hline
 Under 68 kg & Over 63 kg \& Not exceeding 59 kg & Under 62 kg & Over 55 kg \& Not exceeding 59 kg \\
 \hline
 Under 73 kg & Over 68 kg \& Not exceeding 63 kg & Under 67 kg & Over 69 kg \& Not exceeding 63 kg \\
 \hline
 Under 78 kg & Over 73 kg \& Not exceeding 68 kg & Under 73 kg & Over 63 kg \& Not exceeding 68 kg \\
 \hline
 Over 78 kg & Over 78 kg \& Not exceeding 58 kg & Over 68 kg & Over 68 kg \\
 \hline
\end{tabular}
\end{center}
\end{samepage}

\item
\input{5.4.txt}

\begin{center}
\begin{tabular}{ |l|l|l|l| } 
 \hline
 \multicolumn{2}{|l|}{Men's division} & \multicolumn{2}{|l|}{Women's division}  \\
 \hline
 Under 48 kg & Not exceeding 48 kg & under 44 kg & Not exceeding 44 kg \\
 \hline
 Under 55 kg & Over 48 kg \& Not exceeding 55 kg & Under 49 kg & Over 44 kg \& Not exceeding 49 kg \\
 \hline
 Under 63 kg & Over 55 kg \& Not exceeding 63 kg & Under 55 kg & Over 49 kg \& Not exceeding 55 kg \\
 \hline
 Under 73 kg & Over 65 kg \& Not exceeding 73 kg & Under 63 kg & Over 55 kg \& Not exceeding 63 kg \\
 \hline
 Over 73 kg & Over 73 kg & Over 63 kg & Over 63 kg \\
 \hline
\end{tabular}
\end{center}

\item
\input{5.5.txt}

\begin{center}
\begin{tabular}{ |l|l|l|l| } 
 \hline
 \multicolumn{2}{|l|}{Men's division} & \multicolumn{2}{|l|}{Women's division}  \\
 \hline
 Under 33 kg & Not exceeding 33 kg & under 29 kg & Not exceeding 29 kg \\
 \hline
 Under 37 kg & Over 33 kg \& Not exceeding 37 kg & Under 33 kg & Over 29 kg \& Not exceeding 33 kg \\
 \hline                                                                                        
 Under 41 kg & Over 37 kg \& Not exceeding 41 kg & Under 37 kg & Over 33 kg \& Not exceeding 37 kg \\
 \hline                                                                                        
 Under 45 kg & Over 41 kg \& Not exceeding 45 kg & Under 41 kg & Over 37 kg \& Not exceeding 41 kg \\
 \hline                                                                                        
 Under 49 kg & Over 45 kg \& Not exceeding 49 kg & Under 44 kg & Over 41 kg \& Not exceeding 44 kg \\
 \hline                                                                                        
 Under 53 kg & Over 49 kg \& Not exceeding 53 kg & Under 47 kg & Over 44 kg \& Not exceeding 47 kg \\
 \hline                                                                                        
 Under 57 kg & Over 53 kg \& Not exceeding 57 kg & Under 51 kg & Over 47 kg \& Not exceeding 51 kg \\
 \hline                                                                                        
 Under 61 kg & Over 57 kg \& Not exceeding 61 kg & Under 55 kg & Over 51 kg \& Not exceeding 55 kg \\
 \hline                                                                                        
 Under 65 kg & Over 61 kg \& Not exceeding 65 kg & Under 59 kg & Over 55 kg \& Not exceeding 59 kg \\
 \hline
 Over 65 kg & Over 65 kg \& Not exceeding 58 kg & Over 59 kg & Over 59 kg \\
 \hline
\end{tabular}
\end{center}

  \begin{enumerate}
  \begin{samepage}
  \item
  \input{5.5.1.txt}

    \begin{center}
    \begin{tabular}{ |l|l|l|l| } 
     \hline
     \multicolumn{4}{|l|}{Men's division} \\
     \hline
     \multicolumn{2}{|l|}{Cadet contestants' Height} & MAX. Weight & MIN. Weight  \\
     \hline
     Under 148 cm & Not exceeding 148 cm & 45 kg & 33 kg \\
     \hline
     Under 152 cm & Over 148 cm \& Not exceeding 152 cm & 48 kg & 35 kg \\
     \hline                                                       
     Under 156 cm & Over 152 cm \& Not exceeding 156 cm & 51 kg & 37 kg \\
     \hline                                                       
     Under 160 cm & Over 156 cm \& Not exceeding 160 cm & 53 kg & 39 kg \\
     \hline                                                       
     Under 164 cm & Over 160 cm \& Not exceeding 164 cm & 56 kg & 41 kg \\
     \hline                                                       
     Under 168 cm & Over 164 cm \& Not exceeding 168 cm & 59 kg & 43 kg \\
     \hline                                                       
     Under 172 cm & Over 168 cm \& Not exceeding 172 cm & 61 kg & 45 kg \\
     \hline                                                       
     Under 176 cm & Over 172 cm \& Not exceeding 176 cm & 64 kg & 47 kg \\
     \hline                                                       
     Under 180 cm & Over 176 cm \& Not exceeding 180 cm & 67 kg & 49 kg \\
     \hline
     Over 180 cm & Over 180 cm & 80 kg & 52 kg  \\
     \hline
    \end{tabular}
    \end{center}
    \end{samepage}

    \begin{center}
    \begin{tabular}{ |l|l|l|l| } 
     \hline
     \multicolumn{4}{|l|}{Women's division} \\
     \hline
     \multicolumn{2}{|l|}{Cadet contestants' Height} & MAX. Weight & MIN. Weight  \\
     \hline
     Under 144 cm & Not exceeding 144 cm & 43 kg & 32 kg \\
     \hline
     Under 148 cm & Over 144 cm \& Not exceeding 148 cm & 45 kg & 33 kg \\
     \hline                                               
     Under 152 cm & Over 148 cm \& Not exceeding 152 cm & 48 kg & 35 kg \\
     \hline                                                       
     Under 156 cm & Over 152 cm \& Not exceeding 156 cm & 51 kg & 37 kg \\
     \hline                                                       
     Under 160 cm & Over 156 cm \& Not exceeding 160 cm & 53 kg & 39 kg \\
     \hline                                                       
     Under 164 cm & Over 160 cm \& Not exceeding 164 cm & 56 kg & 41 kg \\
     \hline                                                       
     Under 168 cm & Over 164 cm \& Not exceeding 168 cm & 59 kg & 43 kg \\
     \hline                                                       
     Under 172 cm & Over 168 cm \& Not exceeding 172 cm & 61 kg & 45 kg \\
     \hline                                                               
     Under 176 cm & Over 172 cm \& Not exceeding 176 cm & 64 kg & 47 kg \\
     \hline                 
     Over 176 cm & Over 176 cm & 75 kg & 50 kg  \\
     \hline
    \end{tabular}
    \end{center}
  \end{enumerate}

\item
\input{5.6.txt}

\begin{center}
\begin{tabular}{ |m{5em}|m{6em}|m{6em}|m{6em}|m{6em}|m{6em}|m{6em}|m{10em}|  } 
 \hline
 Division & Male pair & Female pair & \multicolumn{2}{|c|}{Male Team} & \multicolumn{2}{|c|}{Female Team} & Mixed Team \\
 \hline
 Maximum \newline number of \newline athletes & 2 & 2 & 4 & 3 & 4 & 3 & 4 (Maximum 2 male \& 2 female) \\
 \hline                                                                                        
 \multirow{2}{6em}{Total weight range} & 160 kg or less & 135 kg or less & \multirow{2}{5em}{300 kg or less} & \multirow{2}{5em}{240 kg or less} & \multirow{2}{5em}{260 kg or less} & \multirow{2}{5em}{200 kg or less} & \multirow{2}{10em}{2 female athletes: 135 kg or less 2 male athletes: 160 kg or less} \\
 \cline{2-3}
 & 130 kg or less & 110 kg or less & & & & & \\
 \hline                                                                                        
\end{tabular}
\end{center}

\end{enumerate}


\chapter{Classification and methods of competition}\label{classification-and-methods-of-competition}

\begin{enumerate}

\item
\input{6.1.txt}

  \begin{enumerate}
  \item
  \input{6.1.1.txt}
  \item
  \input{6.1.2.txt}
  \end{enumerate}

\item
\input{6.2.txt}

  \begin{enumerate}
  \item
  \input{6.2.1.txt}
  \item
  \input{6.2.2.txt}
  \end{enumerate}

\item
\input{6.3.txt}
\item
\input{6.4.txt}
\item
\input{6.5.txt}
\end{enumerate}

\begin{expl}
\section{(Interpretation)}\label{interpretation-4}
\input{6.i_1.txt}

\section{(Explanation \#1)}\label{explanation-1-2}
\input{6.e_1.txt}
\end{expl}


\chapter{Duration of Contest}\label{duration-of-contest}

\begin{enumerate}
\item
\input{7.1.txt}

  \begin{enumerate}
  \item
  \input{7.1.1.txt}
  \item
  \input{7.1.2.txt}
  \item
  \input{7.1.3.txt}
  \end{enumerate}
\item
\input{7.2.txt}
\end{enumerate}

\chapter{Drawing of Lots}\label{drawing-of-lots}

\begin{enumerate}
\item
\input{8.1.txt}
\item
\input{8.2.txt}
\item
\input{8.3.txt}
\end{enumerate}

\chapter{Weigh-in}\label{weigh-in}

\begin{enumerate}
\item
\input{9.1.txt}
\item
\input{9.2.txt}

  \begin{enumerate}
  \item
  \input{9.2.1.txt}
  \item
  \input{9.2.2.txt}
  \end{enumerate}

\item
\input{9.3.txt}

  \begin{enumerate}
  \item
  \input{9.3.1.txt}
  \end{enumerate}

\item
\input{9.4.txt}
\item
\input{9.5.txt}
\end{enumerate}

\begin{expl}
\section{(Explanation\#1)}\label{explantion1}
\input{9.e_1.txt}

\section{(Explanation \#2)}\label{explanation-2-2}
\input{9.e_2.txt}

\section{(Explanation \#3)}\label{explanation-3-2}
\input{9.e_3.txt}

\section{(Explanation \#4)}\label{explanation-4-1}
\input{9.e_4.txt}
\end{expl}


\chapter{Procedure of the Contest}\label{procedure-of-the-contest}


\begin{enumerate}
\item
\input{10.1.txt}

\item
\input{10.2.txt}

\item
\input{10.3.txt}

\item
\input{10.4.txt}

  \begin{enumerate}
  \item
  \input{10.4.1.txt}
  \item
  \input{10.4.2.txt}
  \item
  \input{10.4.3.txt}
  \item
  \input{10.4.4.txt}
  \item
  \input{10.4.5.txt}
  \item
  \input{10.4.6.txt}
  \item
  \input{10.4.7.txt}
    \begin{enumerate}
    \item
    \input{10.4.7.1.txt}
    \end{enumerate}
  \item
  \input{10.4.8.txt}
  \end{enumerate}

\item
\input{10.5.txt}
  \begin{enumerate}
  \item
  \input{10.5.1.txt}
  \item
  \input{10.5.2.txt}
  \item
  \input{10.5.3.txt}
  \item
  \input{10.5.4.txt}
  \item
  \input{10.5.5.txt}
  \end{enumerate}

\end{enumerate}

\begin{expl}
\section{(Explanation\#1)}
\input{10.e_1.txt}

\section{(Guideline for officiating)}
\input{10.e_2.txt}
\end{expl}


\chapter{Permitted techniques and areas}\label{permitted-techniques-and-areas}

\begin{enumerate}
\item
\input{11.1.txt}
  \begin{enumerate}
    \item
    \input{11.1.1.txt}
    \item
    \input{11.1.2.txt}
  \end{enumerate}

\item
\input{11.2.txt}
    \begin{enumerate}
    \item
    \input{11.2.1.txt}
    \item
    \input{11.2.2.txt}
    \end{enumerate}
\end{enumerate}


\chapter{Valid Points}\label{valid-points}

\begin{enumerate}
\item
\input{12.1.txt}
    \begin{enumerate}
    \item
    \input{12.1.1.txt}
    \item
    \input{12.1.2.txt}
    \end{enumerate}
\item
\input{12.2.txt}
    \begin{enumerate}
    \item
    \input{12.2.1.txt}
    \item
    \input{12.2.2.txt}
    \item
    \input{12.2.3.txt}
    \item
    \input{12.2.4.txt}
    \end{enumerate}
\item
\input{12.3.txt}
    \begin{enumerate}
    \item
    \input{12.3.1.txt}
    \item
    \input{12.3.2.txt}
    \item
    \input{12.3.3.txt}
    \item
    \input{12.3.4.txt}
    \item
    \input{12.3.5.txt}
    \item
    \input{12.3.6.txt}
    \end{enumerate}
\item
\input{12.4.txt}
    \begin{enumerate}
    \item
    \input{12.4.1.txt}
    \end{enumerate}
\item
\input{12.5.txt}
    \begin{enumerate}
    \item
    \input{12.5.1.txt}
    \end{enumerate}
\end{enumerate}

\begin{expl}
\section{(Explanation\#1)}
\input{12.e_1.txt}

\end{expl}


\chapter{Scoring and publication}\label{scoring-and-publication}

\begin{enumerate}
\item
\input{13.1.txt}
\item
\input{13.2.txt}
\item
\input{13.3.txt}
\item
\input{13.4.txt}
\item
\input{13.5.txt}
\end{enumerate}


\chapter{Prohibited acts and Penalties}\label{prohibited-acts-and-penalties}

\begin{enumerate}
\item
\input{14.1.txt}
\item
\input{14.2.txt}
\item
\input{14.3.txt}
\item
\input{14.4.txt}
    \begin{enumerate}
    \item
    \input{14.4.1.txt}
        \begin{enumerate}
        \item
        \input{14.4.1.1.txt}
        \item
        \input{14.4.1.2.txt}
        \item
        \input{14.4.1.3.txt}
        \item
        \input{14.4.1.4.txt}
        \item
        \input{14.4.1.5.txt}
        \item
        \input{14.4.1.6.txt}
        \item
        \input{14.4.1.7.txt}
        \item
        \input{14.4.1.8.txt}
        \item
        \input{14.4.1.9.txt}
        \item
        \input{14.4.1.10.txt}
        \item
        \input{14.4.1.11.txt}
        \item
        \input{14.4.1.12.txt}
        \item
        \input{14.4.1.13.txt}
        \end{enumerate}
    \item
    \input{14.4.2.txt}
    \end{enumerate}
\item
\input{14.5.txt}
\item
\input{14.6.txt}
\item
\input{14.7.txt}
    \begin{enumerate}
    \item
    \input{14.7.1.txt}
    \end{enumerate}
\item
\input{14.8.txt}
\end{enumerate}

\begin{expl}

\section{(Interpretation)}
\input{14.e_1.txt}


\section{(Explanation\#1)}
%\input{14.e_2.txt}
%todo: vagy ezt a részt én fordítom le, vagy később visszatérek, és filokba teszem a tartalmat

``Gam-jeom''

\begin{enumerate}
\def\labelenumi{\roman{enumi}.}
\item
  Crossing the Boundary Line:\\
   A ``Gam-jeom'' shall be declared when one
  foot of a contestant crosses the Boundary Line. No ``Gam-jeom'' will
  be declared if a contestant crosses the boundary line as a result of a
  prohibited act by the opposing contestant.
\item
  Falling down:\\
 ``Gam-jeom'' shall be declared for falling down.
  However, if a contestant falls down due to the opponent's prohibited
  acts ``Gam-jeom'' penalty shall not be given to the fallen contestant,
  while a penalty shall be given to the opponent. If both contestants
  fall as a result of incidental collision, or in case a contestant who
  received a point with turning kick falls down, no penalty shall be
  given.
\item
  Avoiding or delaying the match:

\begin{enumerate}
\def\labelenumii{\alph{enumii})}
\item
  This act involves stalling with no intention of attacking. A
  contestant who continuously displays a non-engaging style shall be
  given a ``Gam-jeom''. If both contestants remain inactive after three
  (3) seconds, the center referee will signal the ``\,``Gong-gyeok''
  command. A ``Gam-jeom'' will be declared: On both contestants if there
  is no activity from them three (3) seconds after the command was
  given; or on the contestant who moved backwards from the original
  position three (3) seconds after the command was given.
\item
  Turning the back and move away to avoid the opponent's attack should
  be punished as it expresses the lack of a spirit of fair play and may
  cause serious injury. The same penalty should also be given for
  evading the opponent's attack by bending below waist level or
  crouching.
\item
  Retreating from the technical engagement only to avoid the opponent's
  attack and to run out the clock, ``Gam-jeom'' shall be given to the
  passive contestant.
\item
  Pretending injury means exaggerating injury or indicating pain in a
  body part not subjected to a blow for the purpose of demonstrating the
  opponent's actions as a violation, and also exaggerating pain for the
  purpose of elapsing the match time.
\item
  ``Gam-jeom'' shall also be given to the athlete who asks the referee
  to stop the contest in order to adjust the position/fit of protective
  equipment.
\end{enumerate}

\item
  Grabbing or pushing the opponent:

\begin{enumerate}
\def\labelenumii{\alph{enumii})}
\item
  This includes grabbing any part of the opponent's body, uniform or
  protective equipment with the hands. It also Includes the act of
  grabbing the foot or let or hooking the leg with forearm. For pushing,
  it is permitted as a quick impact and a contestant must disengage from
  opponent after one push. The flowing acts shall be penalized.

    \begin{itemize}
    \tightlist
    \item
      Pushing the opponent with prolonged or continuous contact
    \item
      Pushing the opponent out of the boundary line
    \item
      Pushing the opponent in a way that prevents kicking motion or any
      normal execution of attacking movement
    \end{itemize}

\end{enumerate}

\item
  Lifting the leg or cut kick motion shall not be penalized only when it
  is followed by execution of punching or kicking technique in
  combination motion.

\item
  Attacking below the waist: This action applies to an attack on any
  part below the waist. When an attack below the waist is caused by the
  recipient in the course of an exchange of techniques, no penalty will
  be given. This article also applies to strong kicking or stamping
  actions to any part of the thigh, knee or shin for the purpose of
  interfering with the opponent's technique.
\item
  Attacking the opponent after ``Kal-yeo'':

\begin{enumerate}
\def\labelenumii{\alph{enumii})}
\item
  Attacking after Kal-yeo requires that the attack results in actual
contact to the opponent's body.

\item 
  If the attacking motion started before the Kal-yeo, the attack shall not be penalized.

\item 
  In Instant Video Replay, the timing of Kal-yeo shall be defined as the moment that the
  referee's Kal-yeo hand signal was completed (with fully extended arm);
  and the start the attack shall be defined as the moment that the
  attacking foot is fully off the floor.

\item
  If an attack after Kal-yeo did not land on the opponent's body but appeared deliberate and malicious
  the referee may penalize the behavior with a ``Gam-jeom''

\end{enumerate}

\item
  Hitting the opponent's head with the hand: This article includes
  hitting the opponent's head with the hand (fist), wrist, arm, or
  elbow. However, unavoidable actions due to the opponent's carelessness
  such as excessively lowering the head or carelessly turning the body
  cannot be punished by this article.
\item
  Butting or attacking with the knee: This article relates to an
  intentional butting or attacking with the knee when in close proximity
  to the opponent. However, contact with the knee that happens in the
  following situations cannot be punished by this article. When the
  opponent rushes in abruptly at the moment a kick is being executed
  Inadvertently, or as the result of a discrepancy in distance in
  attacking.
\item
  Attacking the fallen opponent: This action is extremely dangerous due
  to the high probability of injury to the opponent. The danger arises
  from the following:

  \begin{itemize}
  \item
    The fallen opponent is in an immediate defenseless state
  \item
    The impact of any technique which strikes a fallen contestant will
    be greater due to the contestant's position. These types of
    aggressive actions toward a fallen opponent are not in accordance
    with the spirit of taekwondo and as such are not appropriate to
    taekwondo competition. In this regard, penalties should be given for
    intentionally attacking the fallen opponent regardless of the degree
    of impact
  \end{itemize}
\end{enumerate}

When misconduct is committed by a contestant or a coach during a rest
period, past the five (5) seconds of the round conclusion, the referee
can immediately declare the ``Gam-jeom'' and the ``Gam-jeom'' shall be
recorded to the upcoming round. However, ``Gam-jeom'' shall be recorded
to the previous round if the action happened within five (5) seconds of
the round conclusion.

\end{expl}


\chapter{Golden Points and Decision of Superiority}\label{golden-points-and-decision-of-superiority}

\begin{enumerate}
\item
\input{15.1.txt}
\item
\input{15.2.txt}
\item
\input{15.3.txt}
\item
\input{15.4.txt}
    \begin{enumerate}
    \item
    \input{15.4.1.txt}
    \item
    \input{15.4.2.txt}
    \item
    \input{15.4.3.txt}
    \item
    \input{15.4.4.txt}
    \item
    \input{15.4.5.txt}
    \end{enumerate}
\item
\input{15.5.txt}
    \begin{enumerate}
    \item
    \input{15.5.1.txt}
    \item
    \input{15.5.2.txt}
    \item
    \input{15.5.3.txt}
    \item
    \input{15.5.4.txt}
    \end{enumerate}
\end{enumerate}

\begin{expl}

\section{(Explanation\#1)}
\input{15.e_1.txt}
\section{(Explanation\#2)}
\input{15.e_2.txt}
\section{(Guideline for officiating)}
\input{15.e_3.txt}
\section{(Guideline for officiating for the best of three (3) system)}
\input{15.e_4.txt}
\end{expl}


\chapter{Decisions}\label{decisions}

\begin{enumerate}
\item
\input{16.1.txt}
\item
\input{16.2.txt}
\item
\input{16.3.txt}
\item
\input{16.4.txt}
\item
\input{16.5.txt}
\item
\input{16.6.txt}
\item
\input{16.7.txt}
\item
\input{16.8.txt}
\item
\input{16.9.txt}
\end{enumerate}

\begin{expl}
\section{(Explanation\#1)}
\input{16.e_1.txt}
\section{(Explanation\#2)}
\input{16.e_2.txt}
\section{(Explanation\#3)}
\input{16.e_3.txt}
\section{(Explanation\#4)}
\input{16.e_4.txt}
\section{(Explanation\#5)}
\input{16.e_5.txt}
\section{(Explanation\#6)}
\input{16.e_6.txt}
\section{(Explanation\#7)}
\input{16.e_7.txt}
\end{expl}

\chapter{Knock Down}\label{knock-down}

A Knock Down shall be declared, when a legitimate attack is delivered and;

\begin{enumerate}
\item
\input{17.1.txt}
\item
\input{17.2.txt}
\item
\input{17.3.txt}
\end{enumerate}

\begin{expl}
\section{(Explanation\#1)}
\input{17.e_1.txt}
\end{expl}


\chapter{Procedure in the event of a Knock Down}\label{knock-down-proc}

\begin{enumerate}
\item
\input{18.1.txt}
    \begin{enumerate}
    \item
    \input{18.1.1.txt}
    \item
    \input{18.1.2.txt}
    \item
    \input{18.1.3.txt}
    \item
    \input{18.1.4.txt}
    \item
    \input{18.1.5.txt}
    \item
    \input{18.1.6.txt}
    \item
    \input{18.1.7.txt}
    \end{enumerate}
\item
\input{18.2.txt}
    \begin{enumerate}
    \item
    \input{18.2.1.txt}
    \item
    \input{18.2.2.txt}
    \item
    \input{18.2.3.txt}
    \item
    \input{18.2.4.txt}
    \item
    \input{18.2.5.txt}
    \item
    \input{18.2.6.txt}
    \item
    \input{18.2.7.txt}
    \end{enumerate}
\end{enumerate}

\begin{expl}
\section{(Explanation\#1)}
\input{18.e_1.txt}
\section{(Guideline for officiating)}
\input{10.e_2.txt}
\section{(Explanation\#2)}
\input{18.e_3.txt}
\section{(Explanation\#3)}
\input{18.e_4.txt}
\section{(Explanation\#4)}
\input{18.e_5.txt}
\section{(Explanation\#5)}
\input{18.e_6.txt}
\section{(Guideline for officiating)}
\input{18.e_7.txt}
\end{expl}


\chapter{Procedures of suspending the match}\label{suspend}

\begin{enumerate}
\item
\input{19.1.txt}
    \begin{enumerate}
    \item
    \input{19.1.1.txt}
    \item
    \input{19.1.2.txt}
        \begin{enumerate}
        \item
        \input{19.1.2.1.txt}
        \item
        \input{19.1.2.2.txt}
        \end{enumerate}
    \item
    \input{19.1.3.txt}
    \item
    \input{19.1.4.txt}
    \item
    \input{19.1.5.txt}
        \begin{enumerate}
        \item
        \input{19.1.5.1.txt}
        \item
        \input{19.1.5.2.txt}
        \end{enumerate}
    \item
    \input{19.1.6.txt}
    \item
    \input{19.1.7.txt}
    \item
    \input{19.1.8.txt}
    \end{enumerate}
\end{enumerate}

\begin{expl}
\section{(Explanation\#1)}
\input{19.e_1.txt}
\section{(Explanation\#2)}
\input{19.e_2.txt}
\end{expl}


\chapter{Technical Officials}\label{officials}

\begin{enumerate}
\item
\input{20.1.txt}
    \begin{enumerate}
    \item
    \input{20.1.1.txt}
    \item
    \input{20.1.2.txt}
    \end{enumerate}
\item
\input{20.2.txt}
    \begin{enumerate}
    \item
    \input{20.2.1.txt}
    \item
    \input{20.2.2.txt}
    \item
    \input{20.2.3.txt}
    \end{enumerate}
\item
\input{20.3.txt}
    \begin{enumerate}
    \item
    \input{20.3.1.txt}
    \item
    \input{20.3.2.txt}
        \begin{enumerate}
        \item
        \input{20.3.2.1.txt}
            \begin{enumerate}
            \item
            \input{20.3.2.1.1.txt}
            \item
            \input{20.3.2.1.2.txt}
            \item
            \input{20.3.2.1.3.txt}
            \item
            \input{20.3.2.1.4.txt}
            \item
            \input{20.3.2.1.5.txt}
            \end{enumerate}
        \item
        \input{20.3.2.2.txt}
            \begin{enumerate}
            \item
            \input{20.3.2.2.1.txt}
            \item
            \input{20.3.2.2.2.txt}
            \end{enumerate}
        \item
        \input{20.3.2.3.txt}
            \begin{enumerate}
            \item
            \input{20.3.2.3.1.txt}
            \end{enumerate}
        \item
        \input{20.3.2.4.txt}
            \begin{enumerate}
            \item
            \input{20.3.2.4.1.txt}
            \item
            \input{20.3.2.4.2.txt}
            \item
            \input{20.3.2.4.3.txt}
            \end{enumerate}
        \end{enumerate}
    \item
    \input{20.3.3.txt}
        \begin{enumerate}
        \item
        \input{20.3.3.1.txt}
        \item
        \input{20.3.3.2.txt}
        \end{enumerate}
    \item
    \input{20.3.4.txt}
        \begin{enumerate}
        \item
        \input{20.3.4.1.txt}
        \item
        \input{20.3.4.2.txt}
        \end{enumerate}
    \item
    \input{20.3.5.txt}
    \item
    \input{20.3.6.txt}
        \begin{enumerate}
        \item
        \input{20.3.6.1.txt}
        \item
        \input{20.3.6.2.txt}
        \end{enumerate}
    \end{enumerate}
\item
\input{20.4.txt}
\end{enumerate}

\begin{expl}
\section{(Explanation\#1)}
\input{20.e_1.txt}
\section{(Interpretation)}
\input{20.e_2.txt}
\section{(Interpretation)}
\input{20.e_3.txt}
\section{(Guideline for officiating)}
\input{20.e_4.txt}
\end{expl}


\chapter{Instant Video Replay}\label{ivr}


\begin{enumerate}
\def\labelenumi{\arabic{enumi}.}
\tightlist
\item
\end{enumerate}

In case there is an objection to a judgment of the refereeing officials
during the contest, the coach of a team can make a request to the center
referee for an immediate review of the video replay. The coach can only
request video replay for followings;

\begin{enumerate}
\def\labelenumi{\roman{enumi})}
\tightlist
\item
  Penalties against the opponent for instances of falling down or
  crossing the boundary line or attacking the opponent after ``Kal-yeo''
  or attacking the fallen opponent
\item
  Technical point
\item
  Any penalty against own contestant
\item
  Any mechanical malfunction or error in time management. In case of
  appeal for PSS mechanical malfunction, the coach may request to the
  center referee for a testing of the PSS at any time during the 2 nd
  and/or the 3rd round. However, if the PSS mechanical function is
  properly working, the coach's appeal quota shall be forfeited. As well
  as this coach's appeal is considered to be a misconduct of coach that
  ``Gam-jeom'' shall be given to own contestant in in accordance with
  Article 14.4.1.3 `Following Misconducts of contestant of coach'. This
  only applies to the `Best of 3' System.
\item
  When referee forgot to invalidate point(s) after ``Gam-jeom'' was
  given for prohibited act
\item
  Wrong identification of fist attacking contestant by judge
\item
  Head kick that is not scored
\end{enumerate}

\begin{enumerate}
\def\labelenumi{\arabic{enumi}.}
\setcounter{enumi}{1}
\tightlist
\item
\end{enumerate}

When coach appeals, the center referee will approach the coach and ask
the reason for the appeal. Any appeal shall not be admissible on any
points scored by foot or fist attacks on the trunk or foot attack on
trunk PSS. Coach may request instant video replay for head kick
regardless of using head PSS. The scope of instant video replay request
is limited to the only one action which has occurred within five (5)
seconds from the moment of the coach's request. Once the coach rises the
blue or red card to request for instant video replay, it will be
considered that the coach has used his/her allocated appeal under any
circumstance unless the judge's meeting satisfies the coach.

3

Referee shall request the Review Jury to review the instant video
replay. Review Jury, who is not of the same nationality as the
contestants, shall review the video replay.

3.1 In the last five (5) seconds of any Round, the Center Referee may
request for IVR review to check the possible Gam-jeom penalties for the
following actions:

\begin{verbatim}
- Falling down
- Crossing the boundary line
- Attack after Kal-yeo
- Attack the fallen opponent
\end{verbatim}

\begin{itemize}
\tightlist
\item
  Any points scored after Prohibited Act will be invalidated.
\end{itemize}

3.2 If a referee perceives a contestant to be staggering, a strong
impact to the head, kick to the eye(s), bleeding or knocked down by a
kick to the head, and so begins counting, but the attack was not scored
by the head PSS, the referee or coach must request IVR to make the
decision for awarding or not awarding points after the count.

3.3 Referee may request IVR for clarification before declaration of
``Gam-jeom'' for pretending injury.

4 After review of the instant video replay, the Review Jury shall inform
the center referee of the final decision within thirty (30) seconds
after receiving the request.

5 Coach shall be allocated with one (1) appeal to request an instant
video replay request per each contest. However, based on the size and
level of the Championships, the Technical Delegate may decide the number
of appeal quota during the head of team meeting. If the appeal is
successful and the contested request is corrected, the coach shall
retain the appeal right for the pertinent contest.

6 The decision of the Review Jury is final; no further appeals during
the contest or protest after the contest will be accepted except errors
in the calculation of scoring in contest of Olympic Qualification
Tournaments, G6 event or above as explained article 21.7.1.

7 In the case that there is a clear erroneous decision from the
refereeing officials and scoring operators on identification of the
contestant or errors in the calculation of scoring as follows

\begin{verbatim}
- Score and Gam-jeom input by Operator
- Misidentification of athletes by Center Referee
\end{verbatim}

any of the refereeing officials shall request for IVR and correct the
decision at any time during the round.

7.1 If errors in calculation of scoring or misidentification of the
athlete were not corrected during the round by the refereeing officials
and the errors subsequently affected the result of the winner of the
contest or the round in Olympic Qualification Tournaments, G6 events or
above, the Technical Delegate (TD) has the right to review the case with
assistance of the Competition Supervisory Board (CSB) to make the
necessary correction within 30 minutes after the pertinent round. If the
errors were identified and confirmed, the TD may request for resumption
of the round from the moment when the error was occurred. In case of
multiple errors, competition will be resumed from the moment when the
first error was occurred.

8 In the case of a successful appeal, the Competition Supervisory Board
may investigate the contest at the end of the competition day and take
disciplinary action against the concerned refereeing officials, if
necessary.

9 In any time during the round, any of the judges can ask for adding or
removing technical points regardless of coach's appeal quota.

10 In the tournament where instant video replay system is not available,
the following protest procedure will be applied.

10.1 In case there is an objection to a referee judgment, an official
delegate of the team must submit an application for re-evaluation of
decision (protest application) together with the non-refundable protest
fee of US\$200 to the Board of Arbitration (Competition Supervisory
Board) within 10 minutes after pertinent contest.

10.2 Deliberation of re-evaluation shall be carried out excluding those
members with the same nationality as that of contestant concerned, and
resolution on deliberation shall be made by majority.

10.3 The members of the Board of Arbitration (Competition Supervisory
Board) may summon the refereeing officials for confirmation of events.

10.4 The resolution made by the Board of Arbitration (Competition
Supervisory Board) will be final and no further means of appeal will be
applied.

10.5 Deliberation procedures are as follows:

10.5.1 A coach or head of team from the protesting nation shall be
permitted to make a brief verbal presentation to the Board of
Arbitration in support of their position. The coach or head of team from
the respondent nation shall be allowed to present a brief rebuttal.

10.5.2 After reviewing the protest application, the contest of the
protest must be arranged according to the criterion of ``Acceptable'' or
``Unacceptable''. 10.5.3 If necessary, the Board can hear opinions from
the referee or judges.

10.5.4 If necessary, the Board can review the material evidence of the
decision, such as the written or visual recorded data.

10.5.5 After deliberation, the Board shall hold the secret ballot to
determine a majority decision.

10.5.6 The Chairperson will make a report documenting the outcome of the
deliberation and shall make this outcome publicly known.

10.5.7 Subsequent process following the decision:

10.5.7.1 Errors in determining the match results, mistakes in
calculating the match score or misidentifying a contestant shall result
in the decision being reversed.

10.5.7.2 Error in application of the rules: When it is determined by the
Board that the referee made a clear error in applying the Competition
Rules, the outcome of the error shall be corrected and the referee shall
be punished.

10.5.7.3 Errors in factual judgment: When the Board decides that there
was a clear error in judging the facts such as impact of striking,
severity of action or conduct, intention, timing of an act in relation
to a declaration or area, the decision shall not be changed and the
officials seen to have made the error shall be reprimanded.


\chapter{Deaf-Taekwondo}\label{deaf-taekwondo}


This article outlines the modifications to the Competition Rules used
for Deaf-Taekwondo. For matters not covered by Article 22 the WT
Competition Rules will apply.

1 Qualification of athlete Contestant must have gone through
classification procedures as outlined in the World Para-Taekwondo and
DeafTaekwondo Classification Code and been assigned Sport Class and
Sport Class Status

2 Weight Categories Olympic weight categories apply to competitions in
Deaf-Taekwondo

3 World Deaf-Taekwondo Championships will be organized based on the most
recent Standing Procedure of the World Deaf-Taekwondo Championships.


\chapter{Sanctions}\label{sanctions}

1 The WT President, Secretary General, or Technical Delegate may request
that an on-spot Extraordinary Sanctions Committee be convened for
deliberation when inappropriate behaviors may have been committed by a
coach, a contestant, an official, and/or any member of a Member National
Association.

2 The Extraordinary Sanctions Committee shall investigate the matter,
and summon person(s) concerned for confirmation of events.

3 The Extraordinary Sanctions Committee shall deliberate the matter and
determine if disciplinary actions are to be imposed. The result of
deliberation shall be immediately announced to the public. If there is a
finding of violation, a written decision, including the relevant facts,
rules, supporting evidence (such as witness statements), the sanction
imposed, and rationale, shall be given to the sanctioned party as soon
as reasonably practicable, and a copy shall be included in the Technical
Delegate's report.

3.1 Potential violations on Conduct of a Contestant;

3.1.1 Refusing the referee's command to complete the ending procedures
of the match, including not participating in the declaration of the
winner.

3.1.2 Throwing his/her belongings (head protector, gloves, etc.) as an
expression of dissatisfaction with decision.

3.1.3 Not leaving the competition area after the end of a match

3.1.4 Not returning to a match after a referee's repeated command

3.1.5 Not complying with the competition official's ruling or command

3.1.6 Not complying with the Competition Management Officials reasonable
instructions related to the orderly management of the event

3.1.7 Manipulation of scoring equipment, sensors or/and any part of a
PSS

3.1.8 Any serious unsportsmanlike behavior during a match or aggressive
misconduct toward competition officials.

3.2 Potential violations on conduct of a coach, team official, or other
members of a Member National Association;

3.2.1 Complaining about or/and arguing against an official's decision
during or after a round.

3.2.2 Arguing with the referee or other official(s)

3.2.3 Violent behavior or remark toward officials, opponents or the
opposing side, or spectators during a match

3.2.4 Provoking spectators or spreading false rumor

3.2.5 Instructing athlete(s) to participate in misconduct, such as
remaining in the competition area after a match

3.2.6 Violent behaviors such as throwing or kicking personal
belonging(s) or competition material(s).

3.2.7 Not following instructions of competition officials to leave the
Field of Play or Venue

3.2.8 Any other serious misconducts toward competition officials

3.2.9 Any attempt to bribe competition officials

4 Disciplinary actions: Disciplinary actions issued by the Extraordinary
Sanction Committee may vary according to the degree of the violation.
The following sanctions may be given:

4.1 Disqualification of the athlete

4.2 Warning and order to issue official apology

4.3 Removal of accreditation

4.4 Ban from the competition venue

\begin{enumerate}
\def\labelenumi{\roman{enumi})}
\item
  Ban for the day
\item
  Ban for the duration of the Championships.
\end{enumerate}

4.5 Cancellation of result

\begin{enumerate}
\def\labelenumi{\roman{enumi})}
\item
  Cancellation of the match result and all related merits
\item
  Cancellation of WT Ranking points
\end{enumerate}

4.6 Monetary fine of between \$100-to-\$5,000 US dollars per violation.

5 The Extraordinary Sanctions Committee may recommend to the WT or WT at
its own initiative may investigate and determine, that additional
disciplinary actions be taken against the members/MNA involved,
including but not limited to longer-term suspension, lifetime ban,
and/or additional monetary fines. Such recommendation can be based on
violations of the Competition Rules and Interpretations as well as
violations of the WT Code of Ethics or other pertinent WT rules.

\chapter{Other matters not specified in Competition Rules}\label{other}

\begin{enumerate}
\item
\input{24.1.txt}
    \begin{enumerate}
    \item
    \input{24.1.1.txt}
    \item
    \input{24.1.2.txt}
    \end{enumerate}
\end{enumerate}

\end{document}
